\documentclass[12pt]{article}
\usepackage[utf8]{inputenc}
\usepackage[T1]{fontenc}
\usepackage{hyperref}
\title{Pray 3 Language Resources for K-12 Curriculum Development}
\author{Compiled by Claude Sonnet 5.5 (claude-sonnet-5-5) for Andr\'as Kornai}
\date{}
\begin{document}
\maketitle

\section{Introduction}
Pray 3 (ISO 639-3 \texttt{pry}) is spoken by the Pray people of Nan Province, northern Thailand. The name ``Pray 3'' is used by social scientists to distinguish this group from the related Phai (``Pray 1''). Sources classify it as Austroasiatic (Khmuic, Mal--Phrai); the family label ``Bookkeeping'' in the task identification does not match the sources found. Speaker counts vary widely between sources (from about 6,000 to about 40,000). The searches found no school materials, primers or dictionaries in the language, and this document does not claim any exist.

\section{Primary Educational Resources}
No primers, textbooks or school programmes for Pray 3 turned up in the searches.

\section{Linguistic and Reference Materials}
\begin{itemize}
\item Pray grammar at Ban Pae Klang, Thung Chang district, Nan Province (MA thesis, Mahidol University, 1989, Mingkwan Malapol), open access: \url{https://repository.li.mahidol.ac.th/items/c69617e4-bcb5-466b-8867-98c34a037116/full}
\item Prai language (encyclopedic overview and classification): \url{https://en.wikipedia.org/wiki/Prai_language}
\item Verbix endangered-languages entry for Prai: \url{https://docs.verbix.com/EndangeredLanguages/Prai}
\item ASJP word list entry PRAI 2: \url{https://asjp.clld.org/languages/PRAI_2}
\end{itemize}
Secondary sources mentioned in search summaries, but not located online: a sociolinguistic survey of the Prai of northern Thailand (Jordan-Diller and Diller, 2004) and a comparative study of Pray and Mal phonology (Singnoi, 1988).

\section{Religious and Community Resources}
\begin{itemize}
\item Pray people (ethnographic overview): \url{https://en.wikipedia.org/wiki/Pray_people}
\item Joshua Project people-group profile: \url{https://joshuaproject.net/pray/cards/ROG3/TH/Ctry/PeopNameInCountry}
\end{itemize}

\section{Digital and Audio Resources}
No audio, video or corpus resources were found.

\section{Curriculum Development Context}
Little usable material exists. A curriculum would have to begin with documentation: the 1989 grammar as a starting point, community-led orthography work, and word lists. Instruction would presumably also rely on Thai, the national school language.
\end{document}